\chapter{Національна програма інформатизації}

Мета Національної програми інформатизації (НПІ) полягає у створенні єдиного цифрового простору як проміжного етапу розбудови зрілого інформаційного суспільства. Це прозоре правове середовище, яке надійно захищене, базується на передових міжнародних стандартах, повноцінно забезпечує інформаційні потреби та реалізує права громадян на доступ до достовірної інформації, істотно підвищуючи при цьому ефективність державного управління.

Національна програма інформатизації бере свій початок від Закону України №~74/98-ВР 1998 року та продовжує розвиток у чинному Законі 2024 року №~2807-IX. Головним чином НПІ закріплює протоколи ініціації, контролю та завершення програм інформатизації як на державному, так і на місцевому рівнях. Вона чітко визначає ключових суб'єктів: генерального замовника (Міністерство цифрової трансформації), керівників проєктів, виконавців та спеціалізованих підрядників. Крім того, Програма регламентує життєвий цикл розробки програмного та апаратного забезпечення згідно зі стандартами ISO/IEC~12207 та ISO~9001, а процеси супроводу~--- згідно з ISO/IEC/IEEE~14764:2022.

Основне завдання НПІ~--- безперервна оптимізація підприємств та установ у структурі органів виконавчої влади (ОВВ). Для досягнення цієї мети система державного управління повинна мати дієві інструменти інституційної самотрансформації та забезпечувати узгоджену архітектуру ІТ-комплексів.

\section{Структурне ядро ОВВ}

На мета-рівні безперервний процес реформування органів виконавчої влади (ОВВ) спирається на три фундаментальні складові. У структурі кожного міністерства як структурне ядро сучасного цифрового відомства мають бути наявні:

\begin{enumerate}
    \item \textbf{Управління юстиції та внутрішнього контролю}~--- необхідний атрибут для забезпечення незалежного внутрішнього аудиту, комплаєнсу та службових розслідувань.
    \item \textbf{Департамент інформаційних систем (ІТ-департамент)}~--- виступає як окрема юридична особа з власним кодом ЄДРПОУ, що виконує роль продуктового власника (Product Owner) для цифрових рішень і відповідає за цілісність державних реєстрів.
    \item \textbf{Управління інституційної трансформації}~--- підрозділ, що автоматизує та оптимізує організаційні бізнес-процеси за допомогою впровадження ІТ-продуктів.
\end{enumerate}

Решта профільних управлінь додаються залежно від специфіки міністерства, але ці три стовпи разом із патронатною службою є обов'язковим фундаментом. Правоохоронні та контролюючі органи (НАЗК, НАБУ, Мін'юст тощо) повинні мати безпосередній або опосередкований доступ до першого компонента (управління юстиції), що суворо регламентується правилами ABAC (Attribute-Based Access Control).

Зі свого боку, Мінцифра, як координатор усього метапроцесу цифрової трансформації, отримує авторизований доступ до третього компонента (управління трансформації). Архітектура та шини обміну даними між усіма системами електронного документообігу (СЕД) і реєстрами координуються ІТ-управлінням Мінцифри та екосистемою «Дія».

\section{Стратегічні завдання цифровізації}

Завдання державної цифровізації та інституційної трансформації передбачає виконання таких стратегічних цілей:
\begin{itemize}
    \item кожне міністерство повинно мати свій автономний і висококомпетентний ІТ-департамент, який операційно обслуговує відповідні державні реєстри;
    \item міністерства створюють управління трансформації, що працюють на базі єдиної платформи для зручного моделювання бізнес- та адміністративних процесів (BPMN);
    \item Мінцифра розробляє та затверджує єдині регламенти роботи для ІТ-управлінь та управлінь трансформації у масштабах уряду.
\end{itemize}

Платформа «Дія» при цьому забезпечує не лише базову шину даних та коробкове рішення «Дія: Документообіг», але й виступає провайдером ліцензій для інших учасників ринку (наразі видано понад 30 таких ліцензій).

Оскільки соціоінформаційні системи повинні залишатися максимально автономними та мати керований життєвий цикл, найефективніше реалізовувати ІТ-департаменти у формі окремих державних підприємств. Цей підхід убезпечує ключові інтелектуальні ресурси від ручного політичного втручання та гарантує стабільний захист персональних даних громадян.

\section{Модель міжвідомчої взаємодії}

Для архітектурного проєктування міжвідомчої взаємодії розглядається базова перелікова модель ОВВ, що формує каркас державної інформаційної системи:

\begin{enumerate}
    \item Кабінет Міністрів України;
    \item Міністерство освіти і науки України;
    \item Міністерство охорони здоров'я України;
    \item Міністерство оборони України;
    \item Міністерство соціальної політики України;
    \item Міністерство у справах ветеранів України;
    \item Міністерство закордонних справ України;
    \item Міністерство юстиції України;
    \item Міністерство внутрішніх справ України;
    \item Міністерство економіки України;
    \item Міністерство енергетики України;
    \item Міністерство фінансів України;
    \item Міністерство розвитку громад і територій України;
    \item Міністерство цифрової трансформації України;
    \item Міністерство захисту довкілля та природних ресурсів України.
\end{enumerate}

\newpage
\section{Протокол державної інституційної трансформації}

Для ефективного управління структурними змінами в режимі реального часу пропонується впровадити протокол \textbf{«ОВВ 2.0»}. Він є розширенням базового протоколу системи електронного документообігу (згідно з постановою КМУ №55). Його формальним втіленням стає протокол \textbf{ДІТ} (Державна інституційна трансформація).

Цей спеціалізований інформаційний протокол визначає такі базові операції над організаційною структурою відомства та її політиками:

\begin{enumerate}
    \item Створення, злиття та ліквідація структурних підрозділів;
    \item Розробка, погодження та модифікація установчих і конституційних документів;
    \item Проєктування технічних завдань для розгортання інформаційних систем підрозділів;
    \item Збір вимог, вибір та подальше впровадження відповідних програмних продуктів;
    \item Запуск, супровід та моніторинг операційної діяльності структурних підрозділів.
\end{enumerate}

У процесі функціонування як класичного операційного документообігу, так і інституційного протоколу ДІТ формуються криптографічні зліпки кваліфікованих електронних підписів (КЕП) відповідальних посадових осіб. Ці незмінні сліди дозволяють здійснювати надійний аудит системи~--- вони перевіряються внутрішніми контролюючими органами (управліннями юстиції ОВВ) та безпосередньо координуючим органом~--- Міністерством юстиції України.
